\chapter{epoch围栏、干净选举、分叉修复与重新入组}
本章将副本状态变化串成受控故障切换。这里必须牢记模型限制：本课程的 authority 是单 JVM 内可信的 metadata/HW 控制面，leader 切换由课程显式触发，不是 KRaft 控制器共识，也不覆盖网络分区下的双主仲裁。follower 通过真实 TCP 读取 leader 数据，每个 broker 仍写自己的独立日志。课程选择干净候选，不以数据丢失换取可用性；没有候选时服务必须保持不可用。

\begin{figure}[ht]
\centering
\begin{tikzpicture}[font=\small,node distance=12mm]
\node[draw,rounded corners,minimum width=26mm] (stop) {leader A offline};
\node[draw,rounded corners,minimum width=28mm,right=of stop] (elect) {elect B, epoch+1};
\node[draw,rounded corners,minimum width=31mm,right=of elect] (repair) {reconcile log tail};
\node[draw,rounded corners,minimum width=27mm,right=of repair] (admit) {admit replica};
\node[draw,rounded corners,minimum width=27mm,below=of admit] (write) {ALL produce};
\draw[-{Stealth},thick] (stop)--(elect); \draw[-{Stealth},thick] (elect)--(repair);
\draw[-{Stealth},thick] (repair)--(admit); \draw[-{Stealth},thick] (admit)--(write);
\node[below=3mm of repair,align=center] {旧 epoch 被拒\\HW 以下不可改};
\end{tikzpicture}
\caption{角色切换必须先围栏旧角色，再修复并证明副本完整，最后才可重新进入 ISR。}
\end{figure}

\lessonsection{25}{干净选举与epoch}
\goals{旧 leader 离线时只选择保有提交前缀的在线副本，并使旧角色请求失效。}{理解选举候选资格、reported LEO、HW 下界与 epoch 围栏各自解决的问题。}
\background{先修\stepref{24}{并发复制broker}，并回看\stepref{22}{ISR 与 HW}。leader 是当前接受写入的 broker，follower 保存自己的分区日志。leader 故障后，选举决定由哪份已知日志继续服务；不是随便挑一个在线节点。leader epoch 是角色任期编号，角色改变时递增；请求携带 epoch，broker 用它识别旧任期。epoch 能拒绝旧 broker 的后续请求，但不能回滚它先前产生的磁盘或业务副作用。本课程 authority 是可信单 JVM 控制面，手动触发干净选举；无安全候选时宁可不可用，不做 unclean election。}
\providededit{provided 的 \pathcode{ClusterAuthority}（\pathcode{provided/src/main/java/io/simplekafka/cluster/ClusterAuthority.java}）提供 \method{public boolean isOnline(int brokerId)}、\method{public PartitionMetadata metadata(TopicPartition tp)} 与 \method{public ReplicationSnapshot snapshot(TopicPartition tp)}；\pathcode{ReplicationSnapshot} 位于 \pathcode{provided/src/main/java/io/simplekafka/cluster/ReplicationSnapshot.java}，用于核对 ISR/HW/reportedLEO。三个 broker 的网络与本地日志已提供。}{\pathcode{exercises/src/main/java/io/simplekafka/replication/LeaderElection.java}：\method{public PartitionMetadata elect(TopicPartition tp)}；\method{public void checkLeader(int brokerId,TopicPartition tp,int epoch)}。}
\paragraph{输入、结果与状态。}\pathcode{elect} 读取一个已知 topic-partition 的 leader、epoch、旧 ISR、在线状态、reportedLEO 与 HW；无单独的传输参数。leader 仍在线则原样返回 metadata，不做状态变更。成功切换只变 authority 角色/进度，不直接改任何 broker 的磁盘。 \pathcode{checkLeader} 验证 broker id、分区及请求 epoch。
\lessoncontract{候选必须同时属于选举开始时的旧 ISR、当前在线，且 authority 已报告的 LEO\(\ge HW\)；物理日志看起来更长不能替代报告。符合资格者中选 broker id 最小者，不按 LEO 最大值选。成功时 epoch 恰加一、HW 不回退、ISR 重置为仅新 leader，其他 replica 的 authority reportedLEO 重置为 0 并进入恢复；它们必须修复再 admission 才可参与确认。旧 epoch 请求报 \method{FENCED_EPOCH}；epoch 正确但 broker 非在线 leader 报 \method{NOT_LEADER}。旧 leader 的 ACK waiter 也须被唤醒。}
\begin{itemize}
\item 已知 leader 在线：返回现 metadata，epoch/ISR/HW/进度不变。没有候选：\method{NO_ELIGIBLE_LEADER}，metadata、epoch、HW、ISR 不变；在线但不属于旧 ISR 的 broker 不能补位。
\item 负 broker id/epoch 为 \method{INVALID_REQUEST}；不存在分区为 \method{UNKNOWN_TOPIC_OR_PARTITION}。选举与检查均不发 TCP RPC。
\end{itemize}
\lessonhints{列候选时依次应用旧 ISR、online 与 reportedLEO\(\ge HW\)，不要先按“看起来最新”排序。}{把旧 epoch 与错误角色区分：前者是 \method{FENCED_EPOCH}，后者是 \method{NOT_LEADER}。}{分别试 leader 仍在线、HW 以下 ISR 副本、在线但非 ISR 副本，以及完全无候选；比较调用前后的 authority snapshot。}
\expected{手算前提：leader 1 离线，epoch=0，旧 ISR={1,2,3}，HW=3；日志从 offset 0 起无 retention，三条共同前缀记录 key=null、每个 value 为一个 ASCII 字节、timestamp 固定，故每条33 bytes。副本2 的 reportedLEO=3，副本3 的 reportedLEO=5，二者在线且都属于 ISR。两者都满足 HW 下界，故选 id 较小的 broker 2，即使 broker 3 更长；结果 epoch=1、leader=2、ISR={2}、HW 仍为3。其他 broker 的 authority LEO 被重置为0，但选举不删除其物理日志。随后 broker 3 带 epoch 1 请求会得 \method{NOT_LEADER}；带 epoch 0 的旧请求得 \method{FENCED_EPOCH}。若候选都低于HW或离线，即使 broker 4 在线也必须报告无候选且不改状态。}
\paragraph{测试与完成标准。}\method{Step25Test.electsTheLowestEligibleIsrAndFencesOldEpochWithoutLosingCommittedPrefix} 检查在 leader 在线时重复选举无变化、停机切换、HW 保持及旧 epoch 围栏；\method{electionRequiresTheCommittedPrefixBeforeChoosingTheLowestIsr} 检查 reportedLEO 低于 HW 的较小 id 被跳过；\method{electionFencesAnOutstandingAllAckWithoutRollingBackAppend} 检查待确认 append 不被回滚。单步命令定位 Step25；累计命令从 Step1 到25，单步绿不代表第7章完成。
\pitfalls{只挑在线且 LEO 最大的副本可能引入不在旧 ISR 的分叉；无候选时偷偷选落后副本会丢掉已知提交前缀。真实 Kafka 的 controller、epoch 历史与共识超出课程模型；本课程不保证 authority 丢失或网络分区时的生产级 HA。}
\lessoncommand{25}
\lessonquestions{候选 reportedLEO\(\ge HW\) 能保证什么最小安全条件？}{为何 epoch 正确但不是 leader 的 broker 仍不能接受 produce？}
\paragraph{自检答案。}该条件保证候选至少拥有 authority 已标记为提交的前缀；它不证明未提交尾部也相同。epoch 识别任期，leader 身份决定当前写入权限，两项都必须成立。

\lessonsection{26}{分叉日志修复}
\goals{比较旧副本与新 leader 的相同 offset 内容，保护共同提交前缀，只修复允许变化的尾部。}{理解全前缀扫描、恢复读取、HW 严格边界与 proof 的有效范围。}
\background{先修\stepref{25}{干净选举与 epoch}，并使用\stepref{8}{尾部截断}与\stepref{21}{拉取复制}中的日志/恢复读取概念。LEO 相同只表示条数相同，不能证明内容相同；修复需按 offset 比较每一条 \pathcode{RecordData}。课程用易验证但较慢的全前缀扫描：从0读到捕获的 leader LEO，恢复读取可查看高于HW的 leader 尾部，但不报告复制进度。HW 为排他边界，因此冲突 offset<captured HW 属于已提交内容，拒绝修复；offset 恰好等于 HW 属于未提交尾部，可以修。若存在多个可能的截断位置，截断点不得小于 captured 或最终复核时的当前 HW。}
\providededit{provided 的 \pathcode{RpcClient.call(Messages.Request):Messages.Reply}（\pathcode{provided/src/main/java/io/simplekafka/transport/RpcClient.java}）、\pathcode{Messages.ReplicaFetchRequest/ReplicaFetchBody}（\pathcode{provided/src/main/java/io/simplekafka/protocol/Messages.java}）、\pathcode{RecoveryProof}（\pathcode{provided/src/main/java/io/simplekafka/cluster/RecoveryProof.java}）及 \pathcode{ClusterAuthority} 的 epoch/HW 快照与分区锁已提供；恢复不自动加入 ISR。}{\pathcode{exercises/src/main/java/io/simplekafka/replication/ReplicaReconciler.java}：\method{public long reconcile(int expectedEpoch)}，已提供的 \method{public Optional<RecoveryProof> recoveryProof()}；网络请求为 \pathcode{REPLICA_FETCH(recoveryRead=true)}，每批最多4条、4,200,000 bytes。}
\paragraph{输入、输出与状态。}\pathcode{expectedEpoch} 必须与当前 epoch 相符；对象的 broker id、topic-partition、本地 log、借用 RPC client 与 authority 在构造时确定。reconcile 从 offset0捕获 leader reportedLEO/HW 与本地 mutationVersion，RPC 全在锁外。成功返回本地 end offset，并登记 proof；恢复读取本身不更改 ISR/reportedLEO/HW。进入时清除旧 proof；验证失败不发布新 proof。
\lessoncontract{扫描 leader 捕获前缀与本地完整前缀，找到第一个实际内容差异或尾部结束点。}
\begin{itemize}
\item 预期 epoch 过期为 \method{FENCED_EPOCH}；分区不存在为 \method{UNKNOWN_TOPIC_OR_PARTITION}；本地 broker 是 leader 或未分配为 \method{INVALID_REQUEST}；本地 retained start 不为0、leader前缀不可读/不连续、已提交冲突或截断低于最终HW为 \method{CORRUPT_RECORD}。
\item 本地 offset<captured HW 且与 leader 内容不同，报 \method{CORRUPT_RECORD}，保持文件原样；缺少本地记录（包括缺少提交前缀）可由 leader 补齐。恰好 offset=captured HW 的首个分歧允许改写。仅比 leader 多出的未提交后缀可从 leader LEO 截掉。
\item RPC 超时或本地日志 API 在扫描期间改变 mutationVersion（即使 truncate 后追加使 LEO 回到原值）时报 \method{REQUEST_TIMEOUT}，不修复且不发布 proof。epoch/leader 在扫描期改变报 \method{FENCED_EPOCH}。错误 RPC code 传播；unexpected body 为 \method{INVALID_REQUEST}；本地磁盘错误保留 \method{STORAGE_ERROR}。截断/追加中的I/O不保证跨文件原子回滚。
\item 最终在 authority lock→本地 log monitor 顺序下复核 epoch、leader、HW、start、LEO、mutationVersion，持有日志 monitor 到修复和 proof 发布。proof 绑定捕获 epoch、leader LEO、HW、本地 log 对象身份及最终 mutationVersion；它不是持久校验，也不能检测绕过 API 的外部磁盘篡改。
\end{itemize}
\lessonhints{比较“本地缺记录”与“同 offset 有不同内容”；二者不意味着同一种损坏。}{把HW想成不含该 offset 的排他边界：小于它受保护，等于它可以成为首个可修复位置。}{分别构造共同提交前缀+未提交分歧尾部、短副本/空副本、同LEO并发替换；检查磁盘与proof是否改变。}
\expected{假设 offset 0–2 是共同提交记录，key=null、每条 value 是一个 ASCII 字节、每条33 bytes；HW=3，捕获的 leader LEO=4。旧副本 offset3 为 \pathcode{X}，新 leader offset3 为 \pathcode{Y}，二者 timestamp 相同。offset0–2 相同，首差异恰在3=HW，修复后本地 offset3 为Y、LEO=4并产生 proof；authority ISR/HW不变。另一个独立状态若 offset2 已不同（2<HW），报 \method{CORRUPT_RECORD}，本地文件字节和 authority 不变且旧 proof 已清除。若本地为空则补齐0–3，不属于已提交冲突。}
\paragraph{测试与完成标准。}\method{Step26Test.repairsOnlyDivergentUncommittedTailAtHighWatermarkAndReturnsLeaderPrefixProof} 检查 offset=HW 可修及proof；\method{refusesMismatchAtHighWatermarkMinusOneWithoutChangingTheLogOrAuthorityAndClearsOldProof} 检查 offset=HW-1 的拒绝；\method{fillsShortAndEmptyReplicasWithTheCommittedPrefix}、\method{trimsOnlyAnUncommittedTailWhenTheLocalReplicaHasTheCompleteLeaderPrefix} 与 \method{repairsTenRecordsInFourRecordFetchRoundsWithoutReportingRecoveryProgress} 检查缺失/多余尾部和4条分页；另有 RPC超时、畸形回复、same-LEO mutation、epoch/HW变化竞态。单步命令定位Step26；累计命令运行Step1–26。
\pitfalls{本地 HW=0 不表示可以删掉本地已有的有效前缀；缺失记录不是内容冲突。真实 Kafka 用 leader-epoch 历史更有效地找分歧；本课程的 O(n) 全扫描用于教学，不应视为生产修复算法。}
\lessoncommand{26}
\lessonquestions{为何已提交 offset 内容冲突与本地缺少已提交记录必须区别处理？}{leader 捕获快照后又追加，为什么扫描得到的旧 proof 不能直接用于 admission？}
\paragraph{自检答案。}冲突意味着两边对被保护的已提交值不一致，不能静默覆盖；缺失表示副本落后，若 leader 前缀可用则可补齐。leader LEO改变会使proof没有证明最新尾部，所以之后 admission 必须拒绝并重新扫描。

\lessonsection{27}{元数据路由刷新}
\goals{让客户端通过 metadata 把分区请求送到 leader，同时避免因一次错误自动重放业务操作。}{从零理解 bootstrap、缓存路由、“结果未知”及为什么刷新控制面不等于重试数据面请求。}
\background{先修\stepref{26}{分叉日志修复}，并回看\stepref{13}{批量生产者}与\stepref{23}{生产确认}。客户端通常只知道 bootstrap endpoint——用于发现集群的入口；metadata 回答 topic-partition 当前 leader endpoint 与 epoch。缓存的 leader 是数据请求的目标，authority/metadata 请求是控制面，produce/fetch 是数据面。结果未知表示 broker 可能已完成 append，而回复在网络中丢失或 ACK 等待失败，客户端无法从错误单独判断是否写入。若自动重送普通 produce，可能产生第二条记录。}
\providededit{provided 的 \pathcode{RpcClientFactory.connect(Endpoint):RpcClient}（\pathcode{provided/src/main/java/io/simplekafka/transport/RpcClientFactory.java}）、\pathcode{RpcClient.call(Messages.Request):Messages.Reply}（\pathcode{provided/src/main/java/io/simplekafka/transport/RpcClient.java}）、\pathcode{Endpoint}（\pathcode{provided/src/main/java/io/simplekafka/model/Endpoint.java}）与 \pathcode{MessageCodec.encodeRequest/decodeRequest}（\pathcode{provided/src/main/java/io/simplekafka/protocol/MessageCodec.java}）负责连接/frame；broker metadata 由课程服务端提供。现有 RpcClient 单 broker 构造路径仍可用。}{\pathcode{exercises/src/main/java/io/simplekafka/client/MetadataRouter.java}：构造器 \method{MetadataRouter(List<Endpoint> bootstrap,RpcClientFactory factory)} 要求 bootstrap 非空且元素非null、factory非null；\method{public void refresh(String topic)}、\method{public Messages.Reply call(TopicPartition tp,Messages.Request request)}。Router 缓存并拥有 factory 创建的 RpcClient，\method{close()} 关闭这些 client。路由构造器路径也在 exercises client 中：\method{SimpleProducer(MetadataRouter,Partitioner,int)}、\method{SimpleConsumer(MetadataRouter,String)}、\method{GroupConsumer(MetadataRouter,String,String)}；不要移除单 broker RpcClient overload。}
\paragraph{输入、输出与状态。}\pathcode{refresh} 按构造时 bootstrap 列表顺序请求一个 topic 的 metadata；成功缓存该 topic 的 partition metadata 并设置 active bootstrap。连接创建/连接请求失败或 \method{REQUEST_TIMEOUT} 可继续下一个入口；若所有入口失败，传播最后失败，不造假路由。topic 无效或回复含其他 topic 为 \method{INVALID_REQUEST}。非 timeout 的 broker error reply/body 错误立即传播，不尝试下一个入口；任何刷新失败都保留上次成功的缓存。Step27 中 \pathcode{call} 接受 partition 相符的 ProduceRequest / FetchRequest；后续第30步再加入 \pathcode{IDEMPOTENT_PRODUCE}。
\lessoncontract{call 按缓存的 leader endpoint 发出一次业务请求，并携带缓存 epoch。请求类型不支持或请求分区与 \pathcode{tp} 不匹配为 \method{INVALID_REQUEST}；没有 topic/partition metadata 为 \method{UNKNOWN_TOPIC_OR_PARTITION}。}
\begin{itemize}
\item 收到响应 \method{NOT_LEADER} 或 \method{FENCED_EPOCH} 后可 best-effort 刷新缓存，使下一次独立调用用新路由；必须返回本次原始 Reply。即使刷新失败也返回原始 Reply。
\item 其他业务错误、RPC timeout、I/O/storage 异常不自动刷新/重放 produce 或 fetch。特别是调用超时可能发生在 broker append 后，未知结果不能解释为“未发生”。
\item 刷新只是读 metadata；Router 不为当前 call 重试业务操作。后续调用由上层显式决定。关闭后的 router 操作抛 \pathcode{IllegalStateException}。
\end{itemize}
\lessonhints{给一次 \pathcode{call} 只画一个业务 request；刷新 metadata 是另一次控制面请求，不是第二次 produce。}{检查 bootstrap 失败顺序与缓存：哪些连接错误可继续，哪些响应错误应原样终止刷新？}{用旧路由触发 \method{FENCED_EPOCH} 或 \method{NOT_LEADER}，核对返回的原错误、刷新后的缓存和下一次显式请求。}
\expected{手算时假设 bootstrap=[broker1,broker2,broker3]，缓存指向 broker1/epoch0，分区空；每条例子记录为 null key、一个 ASCII value 字节、固定非负 timestamp（33 bytes）。broker1 离线并已选举 broker2/epoch1。旧路由请求收到的原始 \method{NOT_LEADER} 或 \method{FENCED_EPOCH} 回复仍返回给调用方，刷新只把 cache 更新至 broker2/epoch1；该调用不在 broker2 再发produce，因此不会在这里产生receipt。调用方随后显式发一个新请求，若其得到 [0,1)，broker2 只有一条新记录。若请求曾在broker追加后因回复超时，Router也不重放；Step29会观察其不确定结果。}
\paragraph{测试与完成标准。}\method{Step27Test.refreshSkipsUnavailableAndFactoryThrowingBootstrapsBeforeCachingMetadata} 检查 bootstrap 顺序；\method{failedRefreshTimesOutWithoutReplacingCacheAndLaterSuccessfulRefreshReplacesIt} 检查缓存保留；\method{returnsLeaderErrorsUnchangedAndRefreshesOnlyForTheNextRequest} 检查原回复与下一请求路由；\method{routerDoesNotReplayAProduceAfterItsAppendResponseTimesOut}、\method{makesRoutedProducerOutcomeUnknownAfterEveryAppendError} 与 \method{routerDoesNotReplayStorageOrTransportFailures} 检查不重放；另验证借用 router 生命周期和不同 topic 独立路由。单步命令定位Step27；累计命令运行Step1–27。
\pitfalls{把 \method{NOT_LEADER} 当成任何情况下都“绝对没写入”会导致重复；刷新缓存后伪造成功receipt也不正确。普通 \method{SimpleProducer} 不自动重试。真实 Kafka 客户端有更复杂的重试与请求语义；本课程 Router 明确只刷新路由、不重放当前业务请求。}
\lessoncommand{27}
\lessonquestions{为什么刷新 metadata 与重放业务请求必须是两个动作？}{客户端收到错误时，为什么不能只凭错误码断定服务端没有副作用？}
\paragraph{自检答案。}metadata 描述之后请求应去哪里；它不能给先前请求补发结果。超时可能发生在 append 已完成之后，而FENCED也可能来自 ACK 等待阶段；若系统没有针对该请求的成功receipt，调用方就应把结果视为未知而不是凭空判定未写入。

\lessonsection{28}{副本重新入ISR}
\goals{只有新鲜的前缀 proof 才可把已修复副本重新纳入 ISR，并将修复、readmission 与 ALL 可用性连接起来。}{理解 assignment、online/follower 角色、epoch、日志身份和 proof 版本。}
\background{先修\stepref{26}{分叉日志修复}与\stepref{25}{干净选举}。修复把日志变成当前 leader 前缀的候选副本；ISR readmission 是单独的控制面决定，成功后副本才参与 ALL/minISR 和 HW。LEO 相等只说明记录数量相同，不证明记录内容相同。reconciliation 的 proof 把一次前缀比较绑定到 broker、分区、epoch、leader LEO、本地日志对象与 mutationVersion；leader 后续追加、目标日志变更、截断或 retention 都会使证据过期。}
\providededit{provided 的 \pathcode{ClusterAuthority}（\pathcode{provided/src/main/java/io/simplekafka/cluster/ClusterAuthority.java}）管理角色/ISR/进度，\pathcode{RecoveryProof}（\pathcode{provided/src/main/java/io/simplekafka/cluster/RecoveryProof.java}）记录 \pathcode{brokerId,tp,epoch,leaderLEO,localLEO,highWatermark,localLog,localMutationVersion}；副本重开与 authority 状态容器已提供。}{\pathcode{exercises/src/main/java/io/simplekafka/replication/ReplicaAdmission.java}：\method{public boolean tryAdd(int brokerId,TopicPartition tp,int epoch)}。}
\paragraph{输入、输出与状态。}\pathcode{tryAdd} 不做 RPC；输入是目标 broker、分区和预期当前 epoch，具体本地 PartitionLog 在构造时绑定。true 表示已经在 ISR 或刚成功 admission；false 表示当前状态/proof 不允许加入。成功会把目标加进 ISR、以 leader LEO 初始化其 reportedLEO 和 caught-up clock、重算 HW、signal waiters，并消费/清除 proof。
\lessoncontract{新 admission 必须确认目标已分配、在线、是当前 epoch follower；proof 的 broker/tp/epoch 正确且 leaderLEO 等于当前 authority reported leader LEO；本地日志必须是proof绑定的同一对象、mutationVersion相同、logStart=0、localLEO=leaderLEO，且proof HW不高于leaderLEO。authority lock→本地日志 monitor 下验证到 ISR 发布之间不得插入本地日志写。}
\begin{itemize}
\item 负 broker id/epoch 或未分配 broker 为 \method{INVALID_REQUEST}；分区未知为 \method{UNKNOWN_TOPIC_OR_PARTITION}；读取本地日志失败为 \method{STORAGE_ERROR}。
\item epoch过期、目标offline、目标当前为leader，或没有/不匹配的新鲜proof时返回false，不加入ISR。proof内容或日志对象/版本/start/LEO不匹配时还清除proof；必须再次reconcile。任何目标日志 append/delete/truncate/close 或 leader LEO 改变都使旧proof不可用，即使truncate+重写恢复同LEO。
\item 只有目标在线、是 follower 且使用当前 epoch 时，已在 ISR 的重复调用才返回true，不改HW；epoch过期、目标offline或目标是leader时按上一项返回false。首次 admission 后，reportedLEO设为leader LEO、caught-up时间设为当前clock；满足minISR时才可能推进HW，否则仍可成功readmission但ALL预检仍失败。
\end{itemize}
\lessonhints{把proof看成针对特定日志对象和版本的短期证据，而不是可永久复用的布尔值。}{列出 admission 重新核对的当前leader状态，以及目标本地日志的身份/版本/start/LEO。}{先验证无状态变化时true，再让leader追加或目标日志truncate后尝试旧proof，最后重新reconcile并admit。}
\expected{手算前提：分区 offset 0 的第一条普通记录 key=null、value为一个ASCII字节、timestamp=30，占33 bytes；leader1与副本2复制并report到LEO=1，HW=1。leader1停机后broker2成为epoch1 leader、ISR={2}；重启的broker1仍有物理LEO=1，但尚不在ISR。reconcile后proof绑定broker1日志和leader LEO1。若leader2先在offset1追加一条timestamp31的同尺寸记录，当前leader LEO=2；broker1用旧proof调用 \pathcode{tryAdd(1,tp,1)} 返回false，ISR/HW不变，broker1仍LEO1。再次reconcile把offset1补上并产生新proof，再admit返回true，ISR={1,2}且reportedLEO=2、HW可推进到2。}
\paragraph{测试与完成标准。}\method{Step28Test.proofBecomesStaleWhenLeaderAppendsAfterReconciliation} 验证leader追加使proof过期；\method{admissionRequiresCurrentProofAndOnlineAssignedFollowerAndIsIdempotent} 验证无proof/旧epoch/offline/unassigned与重复调用；\method{recoveryProofCannotBeUsedWithAnotherLogObjectWithTheSameRecordsAndVersion} 与 \method{admissionRejectsProofAfterTruncateAndAfterDifferentDataRestoresSameLeo} 验证对象身份及同LEO重写；\method{recoveredProofRestoresISRForAllAcksAndReplicaLogsConverge} 与 minISR 等集成测试验证HW唤醒、ALL和日志收敛。单步命令定位Step28；累计命令运行Step1–28，是第7章完成边界。
\pitfalls{相同LEO不证明前缀相同；只核对proof创建时的epoch而不核对当前leader LEO会漏掉新尾部。真实 Kafka 的ISR准入与复制协议更复杂；课程的authority和proof保证是教学模型，不是生产controller设计。}
\lessoncommand{28}
\lessonquestions{两份日志的LEO相同为什么不能单独证明副本安全？}{为什么成功admission后必须更新进度、重算HW并唤醒ALL waiter？}
\paragraph{自检答案。}同样长度可以包含不同记录；proof比较内容并绑定具体日志版本。副本加入ISR可能使满足minISR的共同最小LEO增加，所以需要重算HW并通知等待线程重新检查确认条件。
